\subsubsection{Saint Venant Restriction}

% \begin{definition}[Classical SV resultants]
% \label{def:sv-resultants}
% Let $\bm s_{\rm SV}(x;\IesanVector)$ denote the classical Saint~Venant stress resultants
% associated with $\hat{\bm u}_{\rm P}(\IesanVector)$.
% With $\IesanStiff$ and $\IesanEvolveSP$ defined earlier,
% \[
%   \bm s_{\rm SV}(x;\IesanVector)=\IesanStiff\,\exp(x\IesanEvolveSP)\,\IesanVector
  
% \]
% \end{definition}

%------------------------------------------------------------
% Class I: properties imposed only on the SV restriction
%------------------------------------------------------------

\begin{definition}[Class I]
\label{def:classI-minimal}
A trace $\mathcal T$ is \textbf{Class I} if its SV restriction satisfies:

\begin{enumerate}[label=(\roman*)]
\item \textbf{(D1) SV strain-exactness.}
The linear map $\IesanVector\mapsto \bm e_{\mathcal T}(0;\IesanVector)$ is invertible.

\item \textbf{(T1) SV strain-affinity.}
There exists a constant matrix $\TraceEvolveEE\in\mathbb R^{6\times 6}$ such that
\[
  \bm e_{\mathcal T}(x;\IesanVector) = (\mathbf 1+x\TraceEvolveEE)\,\bm e_{\mathcal T}(0;\IesanVector)
  \qquad \forall x\in I,\ \forall \IesanVector\in\mathbb R^6.
\]

\item \textbf{(C1) Prismatic SV resultant mapping.}
There exists a constant $\TraceStiffNM\in\mathrm{GL}(6)$ such that
\[
  \bm s_{\rm SV}(x;\IesanVector)=\TraceStiffNM\,\bm e_{\mathcal T}(x;\IesanVector)
  \qquad \forall x\in I,\ \forall \IesanVector\in\mathbb R^6.
\]

\item \textbf{(P1${}^{\rm SV}$) Plane-section normalization on primary SV generators.}
There exists a fixed trace center $\CenterTrace\in\Section$ such that for all
$\IesanVector=(a_\varepsilon,\bm a_\Omega,a_\vartheta,\mathbf 0)$ (i.e. $\bm b_\Section=\mathbf 0$),
the traced strain \emph{at $x=0$} satisfies
\[
  \epsilon = a_\varepsilon+\CenterTrace\cdot \bm a_\Omega,\qquad
  \bm\gamma_\perp=\mathbf 0,\qquad
  \varkappa=a_\vartheta,\qquad
  \bm\kappa_\perp=-\FrameBasis^\times \bm a_\Omega.
\]
\end{enumerate}
\end{definition}

%------------------------------------------------------------
% Matrices as representations of restricted operators
%------------------------------------------------------------

\begin{lemma}[Trace matrix as SV representation]
\label{lem:trace-matrix-exists}
If $\mathcal T$ satisfies (D1), then the SV-restricted linear map
$\IesanVector\mapsto \bm e_{\mathcal T}(0;\IesanVector)$ admits a unique matrix representation
$\TraceDeDpRoot\in\mathrm{GL}(6)$ such that
\[
  \bm e_{\mathcal T}(0;\IesanVector)=\TraceDeDpRoot\,\IesanVector.
\]
This matrix $\TraceDeDpRoot$ is the \emph{trace matrix} of the SV-equivalence class.
\end{lemma}

\begin{proposition}[Class I identities on $\mathcal U_{\rm SV}$]
\label{prop:classI-identities}
Let $\mathcal T$ be Class I. Then, for the SV restriction:

\begin{enumerate}[label=(\alph*)]
\item The prismatic resultant matrix is
\[
  \TraceStiffNM = \IesanStiff\,\TraceDpDeRoot.
\]

\item The strain evolution matrix is
\[
  \TraceEvolveEE = \TraceDeDpRoot\,\IesanEvolveSP\,\TraceDpDeRoot.
\]

\item The SV strain field is
\[
  \bm e_{\mathcal T}(x;\IesanVector)=(\mathbf 1+x\TraceEvolveEE)\TraceDeDpRoot\,\IesanVector
  =\TraceDeDpRoot(\mathbf 1+x\IesanEvolveSP)\IesanVector.
\]

\item The SV derivative matrix is uniquely determined by
\[
  \TraceDeDpOne \;\coloneq\; \partial_{\IesanVector}\bm e_{\mathcal T}'(x;\IesanVector)
  \;=\;\TraceDeDpRoot\,\IesanEvolveSP.
\]
\end{enumerate}
\end{proposition}

\begin{proof}

\begin{enumerate}
\item[(a)] From \eqref{eq:snm0-in-p} and Lem.~\ref{lem:trace-matrix-exists}:
\[
  \snm(0;\IesanVector)=\IesanStiff\,\IesanVector
  =\TraceStiffNM\,\bm e_{\mathcal T}(0;\IesanVector)
  =\TraceStiffNM\,\TraceDeDpRoot\,\IesanVector,
\]
so $\TraceStiffNM=\IesanStiff\,\TraceDpDeRoot$.

\item[(b)] Using SV equilibrium $ \bm s_{\rm SV}'=\mathbb J\,\bm s_{\rm SV}$ and (C1):
\[
  \TraceStiffNM\,\bm e_{\mathcal T}'=\mathbb J\,\TraceStiffNM\,\bm e_{\mathcal T},
\]
so $\bm e_{\mathcal T}'=(\TraceStiffNM^{-1}\mathbb J\TraceStiffNM)\bm e_{\mathcal T}$.
Since $\mathbb J^2=\mathbf 0$, this implies $\bm e_{\mathcal T}$ is affine and
\[
  \bm e_{\mathcal T}'(x;\IesanVector)=(\TraceStiffNM^{-1}\mathbb J\TraceStiffNM)\,\bm e_{\mathcal T}(0;\IesanVector),
\]
hence $\TraceEvolveEE=\TraceStiffNM^{-1}\mathbb J\TraceStiffNM$.
Substituting $\TraceStiffNM=\IesanStiff\TraceDpDeRoot$ and
$\IesanEvolveSP=\IesanStiff^{-1}\mathbb J\,\IesanStiff$ gives
$\TraceEvolveEE=\TraceDeDpRoot\,\IesanEvolveSP\,\TraceDpDeRoot$.

\item[(c)] Integrate $\bm e_{\mathcal T}'=\TraceEvolveEE\,\bm e_{\mathcal T}(0;\IesanVector)$ and
use Lem.~\ref{lem:trace-matrix-exists}.

\item[(d)] Differentiate (c) with respect to $x$ and then $\IesanVector$.
\end{enumerate}
\end{proof}

% \begin{proposition}[Plane-trace matrix block structure]
% \label{prop:plane-trace-structure-rigorous}
% If $\mathcal T$ is Class I, then (P1${}^{\rm SV}$) implies that the trace matrix
% $\TraceDeDpRoot$ has the plane form
% \[
% \TraceDeDpRoot = \begin{pmatrix}
% 1 & \CenterTrace^\top & 0 & \TraceNF^\top \\
% \mathbf{0} & \mathbf{0} & \TraceFT & \TraceFF \\
% 0 & \mathbf{0}^\top & 1 & \TraceTF^\top \\
% \mathbf{0} & -\FrameBasis^\times & \TraceBT & \TraceMF
% \end{pmatrix},
% \]
% with $\FrameBasis\cdot\TraceNF=\FrameBasis\cdot\TraceFT=\FrameBasis\cdot\TraceTF=\FrameBasis\cdot\TraceBT=0$,
% and the remaining blocks $(\TraceNF,\TraceFT,\TraceFF,\TraceTF,\TraceBT,\TraceMF)$
% are free subject only to $\TraceDeDpRoot\in\mathrm{GL}(6)$ (18 parameters).
% \end{proposition}

\subsubsection{Reconstruction}

We use the (mild) abuse of notation $\mathcal T^{-1}$ to denote a chosen
\emph{lifting} (embedding) of one--dimensional beam fields into a three--dimensional
continuum displacement field.  In general this is \emph{not} a true inverse of $\mathcal T$ on
$\CauchySpace$; rather, it is a right--inverse on the class of section--rigid
motions for which the plane--section property (P1) is enforced, and a convenient
parametrization of enriched kinematics when warping amplitudes are present.

\subparagraph{Warping-enriched embedding.}
Fix a set of cross-sectional warping shapes $\bm\Phi(\bm r)\in\mathbb R^{3\times n_w}$
(either finite $n_w$ or an abstract countable basis).
For beam fields
\[
  \bm z(x)=(\bm u(x),\bm\theta(x))\in C^1(I;\mathbb R^6),
  \qquad
  \bm\alpha(x)\in C^1(I;\mathbb R^{n_w}),
\]
define the lifting
\begin{equation*}
\mathcal T^{-1}[\bm z,\bm\alpha](x,\bm r)
\;\coloneq\;
\bm u(x)
+\bm\theta(x)\times\bigl(\bm r-\CenterTrace\bigr)
+\bm\Phi(\bm r)\,\bm\alpha(x).
% \label{eq:def-trace-lifting}
\end{equation*}

% \paragraph{Section-rigid specialization.}
% When $\bm\alpha\equiv\mathbf 0$, \eqref{eq:def-trace-lifting} reduces to the plane--section
% embedding about $\CenterTrace$,
% \[
%   \mathcal T^{-1}[\bm z,\mathbf 0](x,\bm r)
%   =\bm u(x)+\bm\theta(x)\times(\bm r-\CenterTrace).
% \]
% On the restricted class of section--rigid motions for which (P1) is imposed, this
% lifting is a right--inverse:
% \[
%   \mathcal T\big[\mathcal T^{-1}[\bm z,\mathbf 0]\big]=\bm z
%   \qquad \text{(on the P1-normalized section--rigid subclass).}
% \]

\subparagraph{Rigid motions and the $\bm Z_\delta$ map.}
The rigid displacement field 
\(
  \hat{\bm u}_{\rm R}(x,\bm r;\bm q)
  % =\bm q_u+\bm q_\theta\times\bigl(x\FrameBasis+\bm r\bigr),
  % \qquad
  \bm q=(\bm q_u,\bm q_\theta)\in\mathbb R^6.
\) 
admits the section--rigid embedding \eqref{eq:def-trace-lifting} with $\bm\alpha\equiv\mathbf 0$ by setting
\[
  \bm\theta_{\rm R}(x;\bm q)=\bm q_\theta,
  \qquad
  \bm u_{\rm R}(x;\bm q)=\bm q_u+\bm q_\theta\times\bigl(x\FrameBasis+\CenterTrace\bigr),
\]
since then
\[
  \bm u_{\rm R}(x;\bm q)+\bm\theta_{\rm R}(x;\bm q)\times(\bm r-\CenterTrace)
  =\bm q_u+\bm q_\theta\times(x\FrameBasis+\bm r)
  =\hat{\bm u}_{\rm R}(x,\bm r;\bm q).
\]
By the plane--section property (P1), the trace of $\hat{\bm u}_{\rm R}$ is therefore
\[
  \bm z_{\rm R}(x)=
  \begin{pmatrix}
  \bm u_{\rm R}(x;\bm q)\\[2pt]
  \bm\theta_{\rm R}(x;\bm q)
  \end{pmatrix}
  =\bm Z_\delta(x)\,\bm q,
\]
where
\begin{equation}
  \bm Z_{\delta}(\reflenx)
  \coloneq
  \begin{pmatrix}
    \mathbf{1} & \bm{x}_T^{\times}(\reflenx) \\
    \mathbf{0} & \mathbf{1}
  \end{pmatrix},
  \qquad
  \bm{x}_T(\reflenx)\coloneq \reflenx\FrameBasis+\CenterTrace,
  \label{eq:def-z-delta-embed}
\end{equation}
and $\bm Z_\delta(\reflenx)\in \mathrm{GL}(6)$ with inverse obtained by changing the sign of
$\bm{x}_T^{\times}$ in the upper-right block.

\subparagraph{Saint~Venant decomposition in traced variables.}
On the Saint~Venant motion family, write the continuum solution as
\[
  \hat{\bm u}(x,\bm r)
  =\hat{\bm u}_{\rm P}(x,\bm r;\IesanVector)+\hat{\bm u}_{\rm R}(x,\bm r;\bm q),
\]
with $\IesanVector\in\mathbb R^6$ the Saint~Venant parameters and $\bm q\in\mathbb R^6$ the rigid parameters.
Linearity of the trace gives the induced decomposition
\begin{equation}
  \bm z(x)
  =\mathcal T[\hat{\bm u}](x)
  =\TraceDzDpX(x)\,\IesanVector+\bm Z_\delta(x)\,\bm q,
  \label{eq:trace-z-embed}
\end{equation}
where $\TraceDzDpX(x)$ is the \emph{particular trace map} for the Saint~Venant class.

Finally, for a Type~3 formulation with $\bm e(x)=\mathbb B^*\bm z(x)$ and a Class~I trace,
the affine strain structure \eqref{eq:def-T0} implies that $\TraceDzDpX(x)$ is characterized by the linear ODE obtained by differentiating
$\bm e=\mathbb B^*\bm z$ and using $\mathbb B^*\bm Z_\delta\equiv \mathbf 0$:
\[
  \TraceDzDpX' + \mathbb J^\top \TraceDzDpX = \mathbf E_{\rm p}(x),
  \qquad
  \TraceDzDpX(0)=\mathbf 0,
\]
with $\mathbf E_{\rm p}(x)$ determined by the Class~I trace data  $\TraceDeDpRoot$ and $\TraceEvolveEE$, recovering the closed-form expression used previously.
